\section{Direct Sums and Free Modules}

\begin{definition}[Direct Sum of Modules]
    \label{def:direct-sum-modules}%
    Let \(M_1, \ldots, M_k\) be \(R\)-modules. The \emph{direct sum}
    \[
        M_1 \oplus \cdots \oplus M_k
    \]
    is the additive group \(M_1 \times \cdots \times M_k\), with the \(R\)-action given by coordinate-wise scaling
    \[
        r\para{m_1, \ldots, m_k} = \para{r m_1, \ldots, r m_k}.
    \]
\end{definition}

\begin{definition}[Linear Dependence]
    \label{def:linear-dependence}%
    Let \(m_1, \ldots, m_k \in M\). We say the set \(\brce{m_1, \ldots, m_k}\) is \emph{linearly independent}, if and only if
    \[
        \sum_{i = 1}^{k} r_i m_i = 0
    \]
    implies \(r_i = 0\) for all \(i\). Otherwise, we say the set is \emph{linearly dependent}.
\end{definition}

\begin{definition}[Free Generation]
    \label{def:free-generation}%
    A subset \(S \subseteq M\) \emph{freely generates} \(M\), if
    \begin{itemize}
        \item \(S\) generates \(M\); and
        \item any set function \(\functiondomain{f}{S}{N}\) to another \(R\)-module \(N\) extends to an \(R\)-module homomorphism \(\functiondomain{\phi_f}{M}{N}\), such that \(\phi_f\para{s} = f\para{s}\) for all \(s \in S\).
    \end{itemize}

    A module \(M\) is \emph{free}, if it is freely generated by some \(S \subseteq M\). We call \(S\) a \emph{basis}.
\end{definition}

Note that free generating sets may \emph{not} exist!

\begin{example}[Module without Free-Generating Set]
    \label{ex:module-without-free-generating-set}%
    Consider the \(\ZZ\)-module \(\ZZ / 2\). If it were freely generated, the only possibility for a basis is \(S = \brce{1}\). But this is \emph{not} free.

    For example, if \(N = \ZZ\), \(f\para{1} = 1\), we would need a non-zero homomorphism \(\ZZ / 2 \to \ZZ\), but this does not exist.
\end{example}

\begin{remark}
    \(\ZZ^k\) is a free \(\ZZ\)-module -- it behaves `similar' to vector spaces we have seen in linear algebra.
\end{remark}

\begin{proposition}[Basis, Linear Independence, and Unique Linear Combination]
    \label{prop:basis-linear-independence-unique-combination}%
    Let \(S = \brce{m_1, \ldots, m_k} \subseteq M\). Then, the following are equivalent:
    \begin{enumerate}
        \item \(S\) generates \(M\) freely;
        \item \(S\) generates \(M\), and \(S\) is linearly independent;
        \item every \(m \in M\) is \emph{uniquely} expressible as ab \(R\)-linear combination
        \[
            m = \sum_{i = 1}^{k} r_i m_i.
        \]
    \end{enumerate}
\end{proposition}

\begin{proof}
    The fact that 2 is equivalent to 3 is familiar from linear algebra, and the proof is identical.

    For 1 implies 2, let \(S\) generate \(M\) freely, and suppose \(S\) is not linearly independent. Then we can write
    \[
        r_1 m_1 + \cdots + r_k m_k = 0,
    \]
    and up to relabelling, say \(r_1 \neq 0\).

    We define a function \(\functiondomain{f}{S}{R}\), by \(m_1 \to 1\), but \(m_i \to 0\) for all \(i > 1\). This extends to some homomorphism \(\functiondomain{\phi_f}{M}{R}\). We may compute
    \[
        0 = \phi_f\para{0} = \phi_f \para{\sum_{i = 1}^{k} r_i m_i} \neq 0,
    \]
    which gives a contradiction.

    For 3 implies 1, suppose \(S\) expresses every element in \(M\) uniquely, as
    \[
        m = \sum_{i = 1}^{k} r_i m_i,
    \]
    given a set function \(\functiondomain{f}{S}{N}\), define
    \[
        \phi_f\para{m} = \sum_{i = 1}^{k} r_i f\para{m_i}.
    \]

    This map is well-defined by uniqueness of the expression. This is clearly a homomorphism, so \(S\) freely generates \(M\).
\end{proof}

\begin{definition}[Module of Relations, Finitely Presented]
    \label{def:module-of-relations-finitely-presented}%
    Let \(M\) be finitely generated, and let the homomorphism \(\functiondomain{\theta}{R^k}{M}\) surjective give the generators. \(\ker \theta\) is called the \emph{module of relations} for \(\theta\). We say \(M\) is \emph{finitely presented} if \(\ker \theta\) is finitely generated.
\end{definition}

\begin{proposition}[Isomorphism as Modules between \(R^n\) and \(R^m\)]
    \label{prop:isomorphism-modules-rn-rm}%
    Let \(R \neq \brce{0}\) be a ring. Then, if \(R^n \isomorphic R^m\) for some \(n, m \in \NN\), then \(n = m\).
\end{proposition}

\begin{proof}
    We will reduce the case of fields. If \(I \idealsub R\) is an ideal, and \(M\) is an \(R\)-module, consider the submodule \(IM \submodule M\) generated by the set
    \[
        \set{rm}{r \in I, m \in M}.
    \]
    
    We look at the quotient \(M / IM\). If \(r \in I\), then the left-multiplication operation
    \[
        \function{r \times -}{M / IM}{M / IM}{\para{m + IM}}{\para{rm + IM}}
    \]
    is the zero map, since \(rm \in IM\).

    Therefore, we can make \(M / IM\) into an \(R / I\)-module, by
    \[
        \para{r + I} \para{m + IM} = \para{rm + IM}.
    \]

    Now, let \(I \idealsub R\) be maximal.\footnote{This is true for general rings, assuming Axiom of Choice, due to Zorn's Lemma. But if \(R\) is Noetherian, this is automatic.} Then \(R^n / I R^n\) and \(R^m / I R^m\) are both vector spaces on \(R / I\), and hence isomorphic respectively to \(\para{R / I}^n\) and \(\para{R / I}^m\). So \(m = n\) by linear algebra.
\end{proof}